\section*{3 The formulas Ref (author source pp.22:6--22:8)}
Let $F$ be a CNF formula with $n$ variables and $m$ clauses. The key ingredient in the non-automatability result of[2] is expressing by a set of clauses $\Ref(F,s)$ the statement that there is a resolution refutation $D_1,\ldots,D_s$ of length $s$ from the clauses of $F$.
The variables of $\Ref(F,s)$ are $D[u,i,b],V[u,i],I[u,j],L[u,v]$ and $R[u,v]$, where $u,v\in[s],i\in[n],j\in[m]$ and $b\in\bits$. The meaning of $D[u,i,b]$ is that $x_i^b$ appears in $D_u$. The meaning of $V[u,i]$ is that $D_u$ is derived as a weakening of the resolvent of two previous clauses on $x_i$, and the meaning of $I[u,j]$ is that $D_u$ is a weakening of the $j$-th clause of $F$. The meaning of $L[u,v]$ is that the left clause (i.e. that which contains $\neg x_i$) from which $D_u$ was derived is $D_v$, and the meaning of $R[u,w]$ is that the right clause (i.e. that which contains $x_i$) from which $D_u$ was derived is $D_w$. We will also use the variables $V[u,0]$ and $I[u,0]$ to indicate whether $D_u$ is derived from previous clauses or from an initial clause of $F$: in the former case,$I[u,0]$ will be true and $V[u,0]$ false, and in the latter $V[u,0]$ will be true and $I[u,0]$ false. The clauses of $\Ref(F,s)$ encode the following conditions: For each $u,v\in[s],i,i'\in[n],j\in[m]$ and $b\in\bits$,
\begin{align*}
&\exists!k\ V[u,k]\ \&\ \exists!k\ I[u,k]\ \&\ \exists!k\ L[u,k]\ \&\ \exists!k\ R[u,k];\tag{3.1}\\
&V[u,0]\iff\neg I[u,0];\tag{3.2}\\
&\neg L[u,v]\text{ for }v\ge u\ \&\ \neg R[u,v]\text{ for }v\ge u;\tag{3.3}\\
&V[u,i]\ \&\ L[u,v]\implies D[v,i,0];\tag{3.4}\\
&V[u,i]\ \&\ R[u,v]\implies D[v,i,1];\tag{3.5}\\
&V[u,i]\ \&\ L[u,v]\ \&\ D[v,i',b]\ \&\ i\ne i'\implies D[u,i',b];\tag{3.6}\\
&V[u,i]\ \&\ R[u,v]\ \&\ D[v,i',b]\ \&\ i\ne i'\implies D[u,i',b];\tag{3.7}\\
&I[u,j]\ \&\ x_i^b\text{ appears in }C_j\implies D[u,i,b];\tag{3.8}\\
&\neg D[u,i,0]\vee\neg D[u,i,1];\tag{3.9}\\
&\neg D[s,i,b].\tag{3.10}
\end{align*}
It was shown, subsequent to[2], that $\Ref(F,s)$ is hard for resolution whenever $F$ is unsatisfiable[9]. In[2], a variation,$\RRef(F,s)$, is used. $\RRef(F,s)$ expresses the fact that there is a resolution refutation $D_1,\ldots,D_s$ or one contained in $D_1,\ldots,D_s$, from the clauses of $F$. $\RRef(F,s)$ has the same variables as $\Ref(F,s)$ plus a new variable $P[u]$ indicating which of the indices $1,\ldots,s$ are active, i.e. are part of the refutation. The clauses of $\RRef(F,s)$ express the following conditions, which are those of $\Ref(F,s)$ conditioned on the fact that $P[u]$ is true, in addition to three new ones requiring $P[s]$ to be true, and $P[v]$ to be true whenever $P[u]$ and $L[u,v]$ or $R[u,v]$ are true:
\begin{align*}
&P[u]\implies\exists!k\ V[u,k]\ \&\ \exists!k\ I[u,k]\ \&\ \exists!k\ L[u,k]\ \&\ \exists!k\ R[u,k];\tag{3.11}\\
&P[u]\implies(V[u,0]\iff\neg I[u,0]);\tag{3.12}\\
&P[u]\implies\neg L[u,v]\text{ for }v\ge u\ \&\ \neg R[u,v]\text{ for }v\ge u;\tag{3.13}\\
&P[u]\implies(V[u,i]\ \&\ L[u,v]\implies D[v,i,0]);\tag{3.14}\\
&P[u]\implies(V[u,i]\ \&\ R[u,v]\implies D[v,i,1]);\tag{3.15}\\
&P[u]\implies(V[u,i]\ \&\ L[u,v]\ \&\ D[v,i',b]\ \&\ i\ne i'\implies D[u,i',b]);\tag{3.16}\\
&P[u]\implies(V[u,i]\ \&\ R[u,v]\ \&\ D[v,i',b]\ \&\ i\ne i'\implies D[u,i',b]);\tag{3.17}\\
&P[u]\implies(I[u,j]\ \&\ x_i^b\text{ appears in }C_j\implies D[u,i,b]);\tag{3.18}\\
&P[u]\implies(\neg D[u,i,0]\vee\neg D[u,i,1]);\tag{3.19}\\
&P[s]\ \&\ \neg D[s,i,b];\tag{3.20}\\
&(P[u]\ \&\ L[u,v]\implies P[v])\ \&\ (P[u]\ \&\ R[u,v]\implies P[v]).\tag{3.21}
\end{align*}
Notice that giving truth values to the $P[u]$ variables (where $P[s]=1$) reduces $\RRef(F,s)$ to $\Ref(F,s')$ where $s'$ is the number of indices $u$ for which $P[u]=1$.
For an integer $k\ge1$, we define $R^k\Ref(F,s)$ as the formula resulting from substituting each variable $P[u]$ in $\RRef(F,s)$ with the conjunction $\bigwedge_{i=1}^kP_i[u]$ for new variables $P_1[u],\ldots,P_k[u]$. Note that $\RRef(F,s)=R^1\Ref(F,s)$.
Now, let $d,N\ge1$ be integers, and let $x$ be a propositional variable. We associate with $x$ $N^{d-1}\lceil\sqrt N/2\rceil$ new variables $x_{i_1,\ldots,i_d}$, where $i_1,\ldots,i_{d-1}\in[N]$ and $i_d\in[\lceil\sqrt N/2\rceil]$. The fact that we make $i_d$ range over $[\lceil\sqrt N/2\rceil]$ instead of $[N]$ will be important later (specifically in Lemma5.2). The depth-$d$ Sipser functions for $x$ are defined by
\begin{align*}
S_{d,N}^{\wedge}(x)&\stackrel{\rm def}=\bigwedge_{i_1=1}^N\bigvee_{i_2=1}^N\cdots\bigwedge_{i_d=1}^{\lceil\sqrt N/2\rceil}x_{i_1,\ldots,i_d},\\
S_{d,N}^{\vee}(x)&\stackrel{\rm def}=\bigvee_{i_1=1}^N\bigwedge_{i_2=1}^N\cdots\bigvee_{i_d=1}^{\lceil\sqrt N/2\rceil}x_{i_1,\ldots,i_d}
\end{align*}
if $d$ is odd, and
\begin{align*}
S_{d,N}^{\wedge}(x)&\stackrel{\rm def}=\bigwedge_{i_1=1}^N\bigvee_{i_2=1}^N\cdots\bigvee_{i_d=1}^{\lceil\sqrt N/2\rceil}x_{i_1,\ldots,i_d},\\
S_{d,N}^{\vee}(x)&\stackrel{\rm def}=\bigvee_{i_1=1}^N\bigwedge_{i_2=1}^N\cdots\bigwedge_{i_d=1}^{\lceil\sqrt N/2\rceil}x_{i_1,\ldots,i_d}
\end{align*}
if $d$ is even.
We define $\RRef_{d,N}(F,s)$ to be the result of substituting every variable of the form $P[u]$ in $\RRef(F,s)$ with $S_{d,N}^{\wedge}(P[u])$ and every other variable $x$ with $S_{d,N}^{\vee}(x)$. Notice that $\RRef_{d,N}(F,s)$ is a set of depth-$(d+1)$ formulas. But, as we want to prove statements about whether $\RRef_{d,N}(F,s)$ has or does not have small depth-$d$ refutations, we must write it as a set of depth-$d$ formulas. We may do that with only a polynomial increase in size, as the only clauses of non constant size of $\RRef(F,s)$ are those of the form $\neg P[u]\vee\bigvee_i X[u,i]$ corresponding to conditions(3.11), and these clauses will have depth-$d$ after the substitution taking us from $\RRef(F,s)$ to $\RRef_{d,N}(F,s)$. Note that the conversion from $\RRef_{d,N}(F,s)$ written as a set of depth-$d$ formulas to its equivalent set of depth-$(d+1)$ formulas can be carried in tree-like depth-$(d+1)$ LK in linear time. In particular, a tree-like depth-$(d+1)$ LK refutation of the latter set can be turned into a tree-like depth-$(d+1)$ LK refutation of the former set, increasing the size by at most a factor of $N^3$.
